% Basic Notations
\DeclareMathOperator{\st}{\text{s.t.}}

% Set Theory
\DeclareMathOperator{\card}{card}                      % Cardinality

% General Algebra
\DeclareMathOperator{\N}{\mathbb{N}}                    % Set of Natural Numbers
\DeclareMathOperator{\Z}{\mathbb{Z}}                    % Set of Integers
\DeclareMathOperator{\Q}{\mathbb{Q}}                    % Set of Rational Numbers
\DeclareMathOperator{\R}{\mathbb{R}}                    % Set of Real Numbers
\DeclareMathOperator{\C}{\mathbb{C}}                    % Set of Complex Numbers
\DeclareMathOperator{\F}{\mathbb{F}}                    % Number Field (F)
\DeclareMathOperator{\0}{\mathbf{0}}
\DeclareMathOperator{\x}{\mathbf{x}}
\DeclareMathOperator{\spn}{span}
\DeclareMathOperator{\img}{im}                          % Image
\DeclareMathOperator{\Aut}{Aut}                         % Automorphism Group

% Topology
\DeclareMathOperator{\B}{\mathcal{B}}                   % Topological Basis
\DeclareMathOperator{\interior}{int}                     % Interior
\DeclareMathOperator{\dist}{dist}                        % Distance
\DeclareMathOperator{\diam}{diam}                        % Diameter
\DeclareMathOperator{\cl}{cl}                            % Closure

% Complex Analysis
\DeclareMathOperator{\ee}{\mathrm{e}}
\DeclareMathOperator{\Log}{Log}

\newcommand{\ii}{\hyperref[dfn:Complex-Number]{\ensuremath{\mathrm{i}}}}
\newcommand{\barC}{\hyperref[dfn:Extended-Complex]{\ensuremath{\bar{\mathbb{C}}}}}
\newcommand{\Arg}{\hyperref[dfn:Argument]{\ensuremath{\operatorname{Arg}}}}
\newcommand{\Ln}{\hyperref[dfn:Logarithmic]{\ensuremath{\operatorname{Ln}}}}
\newcommand{\Windn}{\hyperref[dfn:Winding-Number]{\ensuremath{n}}}
\renewcommand{\Res}{\hyperref[dfn:Residue]{\ensuremath{\operatorname{Res}}}}
\newcommand{\Holo}{\hyperref[dfn:Holomorphic-Function]{\ensuremath{\mathcal{H}}}}
\newcommand{\Mobius}{\hyperref[dfn:Mobius-Transformation]{\mathfrak{M}}}
\newcommand{\D}{\hyperref[dfn:Unit-Disk]{\mathbb{D}}}
\newcommand{\CTranslation}{\hyperref[dfn:Basic-Mobius-Transformation]{\mathcal{T}_{\mathfrak{M}}}}
\newcommand{\CScaling}{\hyperref[dfn:Basic-Mobius-Transformation]{\mathcal{S}_{\mathfrak{M}}}}
\newcommand{\CRotation}{\hyperref[dfn:Basic-Mobius-Transformation]{\mathcal{R}_{\mathfrak{M}}}}
\newcommand{\CInversion}{\hyperref[dfn:Basic-Mobius-Transformation]{\mathcal{I}_{\mathfrak{M}}}}
\NewDocumentCommand{\XRatio}{mmmm}{\left(#1,#2\hyperref[dfn:Cross-Ratio]{;}#3,#4\right)}

% Hyperref Commands
\newcommand{\Holomorphic}{\hyperref[dfn:Holomorphic-Function]{全纯函数}}
\newcommand{\Analytic}{\hyperref[dfn:Analyticity]{解析函数}}
\newcommand{\Meromorphic}{\hyperref[dfn:Meromorphic-Function]{亚纯函数}}
\newcommand{\AngularForm}{\hyperref[dfn:Angular-Form]{三角表示}}
\newcommand{\PowerSeries}{\hyperref[dfn:Power-Series]{幂级数}}
\newcommand{\UniformConvergent}{\hyperref[dfn:Uniform-Convergence]{一致收敛}}
\newcommand{\ConformalMapping}{\hyperref[dfn:Conformal-Mapping]{共形映照}}
\newcommand{\EssentialSingularity}{\hyperref[dfn:Essential-Singularity]{本性奇点}}
\newcommand{\Taylor}{\hyperref[thm:Taylor]{Taylor 定理}}
\newcommand{\CauchyThm}{\hyperref[thm:Cauchy]{Cauchy积分定理}}
